\par\noindent\SourceLeaf{57}%
\hspace*{\parindent}Altera vero

\[
\begin{aligned}
&n\cos\varphi\bigl(A^{(1)}\sin\varphi+2A^{(2)}\sin2\varphi+.\,.\,.\\&\quad+hA^{(h)}\sin h\varphi+(h+1)A^{(h+1)}\sin(h+1)\varphi+.\,.\,.\bigr),
\end{aligned}
\]

\noindent ut cum priore in unam summam conflari possit, cum in singulis terminis productum ex
$\cos\varphi$ in sinum multipli alicuius ipsius $\varphi$ contineat, ex nota formula

\[
\cos\varphi\sin\lambda\varphi=\frac{1}{2}\sin(\lambda+1)\varphi+\frac{1}{2}\sin(\lambda-1)\varphi,
\]

\noindent mutationem subeat necesse est. Sic in ista membrum, quo $\sin\varphi$ contineatur, ex
termino

\[
n\cos\varphi.2A^{(2)}\sin2\varphi
\]

\noindent procedens, fit $n.A^{(2)}\sin\varphi$. Si generatim ex illa quaesiveris membra, quibus
$\sin(h+1)\varphi$ contineatur, sumendus terminus

\[
hA^{(h)}\sin h\varphi.n\cos\varphi
\]

\noindent (eiusque pars prior

\[
\frac{nhA^{(h)}}{2}\sin(h+1)\varphi)
\]

\noindent cum termino

\[
(h+2)A^{(h+2)}\sin(h+2)\varphi.n\cos\varphi
\]

\noindent (ex hoc autem evoluto pars posterior

\[
\frac{n(h+2)A^{(h+2)}}{2}\sin(h+1)\varphi);
\]

\noindent quo facto evadit ipsius $\sin(h+1)\varphi$ in serie quaesita coefficiens

\[
\frac{nhA^{(h)}}{2}+\frac{n(h+2)A^{(h+2)}}{2}.
\]

$\gamma$) Iam vero, cum series ab una parte ($\alpha$) aequalis sit seriebus ab altera positis
($\beta$), coefficientes homologi, i. e. hoc loco in eiusdem anguli sinum ducti aequales
sint oportet. Sic est ex una parte ipsius $\sin\varphi$ coefficiens

\[
=\nu nA-\nu\frac{nA^{(2)}}{2}
\]

\noindent ex altera autem $A^{(1)}+nA^{(2)}$. Quibus aequatis habebitur:

\[
\nu nA-\nu\frac{nA^{(2)}}{2}=A^{(1)}+nA^{(2)},\qquad\mbox{unde}\qquad A^{(2)}=\frac{2(\nu nA-A^{(1)})}{(\nu+2)n}.
\]

